% 第 6 章　牛顿的三条规则（完整样章）
\chapter{牛顿的三条规则}

\step{反常识开场}

\begin{mainline}
冬季奥运会的冰壶赛道上，运动员把一只约 \(19\) 千克的花岗岩冰壶向前推出，然后松手。从这一刻起，没有任何人再碰它，没有发动机，没有绳子，可它能沿着冰面滑行三十多米，直到撞上另一只壶或者被冰面慢慢磨停。

现在换一个场景：把同样一块石头放在草地上，用差不多的力气推出去。它能滑多远？半米都不到。

同样是“推一下然后松手”，结果差了几十倍。真正让人睡不着觉的还不是草地，而是冰面：松手之后，到底\emph{是谁}在给冰壶“续航”？我们全身的直觉都在说——东西要动起来就得有东西一直推着它；推车的人一停手，车就停，这不是天经地义吗？可冰面偏偏站在直觉的对立面：摩擦越小，物体越“不需要”任何人推，它自己就走得又直又远。如果有一块绝对光滑、无限长的冰面，冰壶会变成什么样？

再把场景推得更远一点。太空舱外，一名宇航员拧螺丝时脱了手，扳手慢悠悠地飘走了。它没有火箭，没有翅膀，周围连空气都没有，可它沿着一条直线一直飞，不加速、不减速、不拐弯，直到撞上什么东西为止。谁在维持它的运动？答案是：没有谁。而这个“没有谁”，正是本章的起点。

我们的日常生活被摩擦和空气包围，以至于一个深刻的规律被埋在了“推车要一直推”的经验下面。这一章要做的，就是把这层厚厚的日常经验刮掉，看看底下那三条干净得近乎刻薄的规则——牛顿三定律。它们管着冰壶、扳手、你坐公交车的每一次前倾，也管着火箭升空。
\end{mainline}

\step{先猜一猜}

在往下读之前，请把这一页当成赌桌，押上你的直觉。不要偷看后面的答案——猜错比猜对更有价值，因为本章要修理的正是直觉。

第一题。想象一块绝对光滑、没有尽头的冰面，光滑到连一丁点摩擦都没有。你把冰壶轻轻推出去，然后松手。接下来它会怎样？（a）慢慢减速，最后停下来；（b）一直匀速直线滑下去，永远不停；（c）越滑越快。请写下你的选择和理由。

第二题。桌上铺一张结实的纸条，纸条上立一只装了半杯水的塑料杯，纸条的一头伸出桌沿。你用尽全力、飞快地沿水平方向把纸条抽出来。杯子会怎样？（a）跟着纸条一起飞出去；（b）原地倒下；（c）几乎纹丝不动。如果改成慢慢地拉纸条呢？

第三题。你站在湖中小船上，用桨使劲把水向后推。水被你推向后，船却向前走。可是你推水多用力，水就推你多用力——如果这两个力一样大、方向相反，它们不该“抵消”吗？船凭什么前进？

先写下你的三个答案。等到本章结尾，我们会回来逐一核对。如果你三题全对，那很好；如果你猜错了，更好——你马上就能亲眼看到自己的直觉是在哪一步拐错了弯。

\step{动手实验}

\begin{experiment}[抽纸条：杯子为什么“懒得动”]
\textbf{材料}：一条结实的纸条（裁一条报纸或厨房纸，宽约五厘米，长约三十厘米）、一只塑料杯、水、一张尽量光滑的桌子。

\textbf{步骤一}：杯子装半杯水，立在纸条上，纸条一端伸出桌沿约十厘米。用手指捏住伸出的一端，\emph{慢慢地、匀速地}水平拉动纸条。观察杯子：它几乎跟着纸条一起走。这就是第二题里“慢拉”的情形，摩擦力有充足的时间拖着杯子走。

\textbf{步骤二}：把杯子放回原位。这一次，捏紧纸条，用最快的速度猛地水平抽出——动作要快、要平，不要向上挑。你会看到纸条“嗖”地出来，杯子只是轻轻一颤，几乎纹丝不动。

\textbf{步骤三}：把杯子装满水（更重），再猛抽一次。杯子更稳了。换一个更轻的物体（比如一只空纸杯）试试，它更容易被带走一点。

\textbf{记录}：慢拉和猛抽，杯子的表现为什么截然不同？重杯子和轻杯子呢？把你的观察写下来，暂时不需要解释。这个实验的全部秘密，将在本章逐步拆开：慢拉与猛抽的区别不在“力的大小”，而在力作用了\emph{多久}；重杯子更“懒”，正是本章要给它的正式名字——惯性。
\end{experiment}

还有一个零成本的观察实验：下次坐公交车或地铁时，站稳扶好，留意车辆\emph{启动}和\emph{刹车}两个瞬间你身体的倾斜方向。连续记录三次。你会发现：刹车时你向前倾，启动时你向后仰。多数人会说“刹车时有股力把我往前甩”。请先记住这个感觉——本章后面要证明，那股“向前的力”根本不存在。

\step{旧模型遇到麻烦}

先替直觉说句公道话。两千多年前，亚里士多德把日常经验总结成一套听起来无懈可击的说法：物体要运动，就得有东西推它；推力一撤，运动就停。这套旧模型解释推车、划船、拉磨全都灵验，灵验到整整两千年里几乎没人觉得需要怀疑它。它就是我们从婴儿时期推玩具车开始一点点攒出来的直觉。

顺便替它分析一下“长寿”的原因：我们恰好住在一个摩擦和空气阻力无处不在的角落里。在这个角落里，任何运动的东西确实都会自己停下来，于是“力撤了运动就停”成了一条百试百灵的经验。旧模型不是愚蠢，它是一份忠实但未经审视的观察报告；它的问题在于把“我们这个角落里的常态”误当成了“宇宙的定律”。要戳破它，就得找到摩擦足够小的场合——而冰面、气垫和太空，正好是这样的场合。

但麻烦一个接一个地来了。

麻烦一：冰壶和太空扳手。松手之后推力已经消失，运动却没有消失。旧模型说“力撤了运动就该停”，可冰壶又滑了三十米。旧模型无法回答：这三十米是谁在推？

麻烦二：伽利略的斜面推理。让小球从斜面滚下，它会加速；让它滚上对面的斜面，它会减速，但几乎能爬回原来的高度。把对面斜面的坡度放缓，小球要滚得更远才停下；坡面越光滑，它爬得越接近原来的高度、走得越远。现在做那个著名的极限动作：把对面的斜面放到完全水平，并且想象表面绝对光滑。小球还爬得回“原来的高度”吗？永远爬不回了——于是它只好沿着水平面一直滚下去，不加速，也不减速。这个推理的锋利之处在于：你不需要真的造出无摩擦的世界，只需要把摩擦一步一步减小，观察趋势指向哪里。趋势清清楚楚：\emph{摩擦越小，物体越不需要任何东西来维持运动}。旧模型要我们相信“力维持运动”，而实验的趋势说，力消失的极限处，运动恰恰最稳定。

麻烦三：方向对不上。竖直上抛一个球，脱手之后它受的力（重力）朝下，可它照样向上飞了好一段才回头。如果力是“维持运动的东西”，向下的力怎么会允许向上的运动继续存在？

三条麻烦指向同一个诊断：旧模型把两件不同的事混成了一件事——它分不清“维持运动”和“改变运动”。冰壶停不下来不是因为有谁维持它，而是因为没有什么去改变它；球向上飞不需要向上的力，只有它速度的\emph{变化}才跟重力有关。要把这两件事分开，我们需要一套全新的分工，需要发明几个今天看来天经地义、当年却石破天惊的概念。

\step{发明新概念}

\begin{mainline}
\textbf{第一个新概念：惯性（牛顿第一定律）。}伽利略的斜面告诉我们：水平方向不受任何推拉的小球，会永远匀速直线滚下去。牛顿把它写成第一条规则：\emph{一个物体，只要所受的外力相互抵消（合力为零），它就保持静止或者匀速直线运动。}物体这种“懒得改变运动状态”的脾气，叫惯性。请注意惯性\emph{不是}一种力：它不是谁在物体背后推着，而只是“运动状态没有外力干预时就不变”这件事本身。第一定律还顺手干了一件大事：它把“匀速直线运动”和“静止”划进了同一类——两者都是“不需要解释的状态”，需要解释的只有\emph{变化}。

\textbf{第二个新概念：力是相互作用，不是库存。}旧模型把力想象成可以装进物体里的东西，像往口袋里装糖：推车的人“把力给了”车，车“带着力”跑，力“用完了”车就停。这套画面全错。新的定义是：\emph{力永远是两个物体之间的相互作用}——一只手推、另一只手被推，缺一不可。它像什么？像握手：握手必须有两只手，你没法跟自己握手，也没法把“一个握手”揣进口袋里备用。它不像什么？它不像水壶里的水，不能“装”在物体里储存，更谈不上“用完”。所以“这个物体含有多少力”是一句病句；物体可以具有速度、具有能量、具有动量（第 9、10 章的主角），但力只是它此刻\emph{正在}与外界发生的相互作用。冰壶松手之后没有任何向前的力作用在它身上——这不但不妨碍它前进，恰恰是它前进得又直又稳的原因。

\textbf{第三个新概念：质量是“对速度改变的抗拒”。}推一辆空购物车，轻轻一推它就窜出去；推一辆装满大米的购物车，同样的力气，它慢吞吞地起步。区别在哪？在于第二个物体更“抗拒”速度的改变。我们给这种抗拒程度起名叫\emph{惯性质量}，简称质量。质量越大，同样的力能造成的速度变化越小。要特别小心：质量不是重量。重量是地球拉这个物体的力（第 8 章细讲），而质量是这个物体“抗拒加速”的内在属性——把购物车和宇航员一起搬到失重的太空站里，重量消失了，可你推那辆装满大米的购物车，它照样慢吞吞，质量一两都没有少。

\textbf{第四条规则：合力决定运动状态如何改变（牛顿第二定律）。}既然运动本身不需要解释，需要解释的就是运动的变化。牛顿说：物体速度改变的快慢（加速度），由它所受的\emph{合力}决定，方向与合力相同；同样的合力下，质量越大，加速度越小。更基本的说法是：\emph{合力决定了物体动量变化的快慢}——这句话本章先记住，第 10 章会把它变成主角。注意公式里写的是“合力”：书放在桌上，重力和支持力都作用在书上，二者抵消，合力为零，所以书不动。书不动不是“因为没有力”，而是因为所有的力加总起来等于零。

\textbf{第五条规则：力永远成双，而且分别落在两个物体上（牛顿第三定律）。}你推墙，墙也推你；桨向后推水，水向前推桨；地球向下拉苹果，苹果也向上拉地球。这一对力大小相等、方向相反，但它们\emph{作用在不同的物体上}，因此永远不能拿来互相“抵消”——抵消只发生在同一个物体所受的力之间。

一旦接受“力成双出现”，生活会突然变成一座第三定律的展览馆。走路：你的脚向后蹬地，地面向前推你——在光滑的冰面上这一步蹬不出去，你立刻就明白向前的那股力是谁给的。游泳：手臂把水向后划，水把你向前送。撞墙手疼：你推墙的力越大，墙推回你的力就越大，疼不是“墙很硬”这么简单，而是墙在一丝不苟地执行第三定律。射击时枪托向后撞肩膀（后座力）：火药燃气向前推子弹，子弹就向后推枪。这些例子看似五花八门，骨架却完全相同——想让自己获得向前的推动，就得把别的东西向后推。
\end{mainline}

现在可以用新概念拆解一个著名的抬杠，它就是本节开头的第三题，也是让无数初学者栽跟头的\textbf{反直觉反例——马与马车悖论}。一匹爱思考的马对牛顿说：“你的第三定律讲，我拉车的力等于车拉我的力，方向相反。那么这两个力加起来永远是零，车怎么可能动起来？我干脆不拉了。”马错在哪？错在把作用在两个不同物体上的力加在了一起。判断“车动不动”，只看\emph{车}受的力：马向前拉车，地面阻力向后拖，前者大，车就加速向前。判断“马动不动”，只看\emph{马}受的力：车向后拽马，而马的蹄子向后蹬地、地面反过来向前推马（静摩擦力），这个向前的推力大于车拽马的力，马就加速向前。那一对相等相反的力从来没有见过面，谈什么抵消？

\begin{center}
\begin{tikzpicture}[scale=1.0]
  % 地面
  \draw[thick] (-4.8,0) -- (4.8,0);
  \fill[pattern=north east lines] (-4.8,-0.28) rectangle (4.8,0);
  % 车（左）与马（右），运动方向向右
  \draw[thick,fill=gray!15] (-3.0,0) rectangle (-1.4,1.0);
  \node at (-2.2,0.5) {车};
  \draw[thick,fill=orange!15] (1.4,0) rectangle (3.0,1.0);
  \node at (2.2,0.5) {马};
  \draw[thick] (-1.4,0.5) -- (1.4,0.5) node[midway,above=1pt] {绳};
  % 车受的力
  \draw[-{Stealth[length=3mm]},very thick,red!70!black] (-2.2,1.25) -- (-0.5,1.25) node[right=2pt] {马拉车（向前）};
  \draw[-{Stealth[length=3mm]},very thick,blue!70!black] (-2.2,0.5) -- (-3.9,0.5) node[left=2pt] {地面阻力};
  % 马受的力
  \draw[-{Stealth[length=3mm]},very thick,red!70!black] (2.2,1.55) -- (0.5,1.55) node[left=2pt] {车拉马（向后）};
  \draw[-{Stealth[length=3mm]},very thick,blue!70!black] (2.2,0.5) -- (3.9,0.5) node[right=2pt] {地面推马};
\end{tikzpicture}
\end{center}

这张图同时引入了本章最后一个工具：\textbf{受力图与系统边界}。做法只有两步：先用一条想象中的虚线把“研究对象”圈出来（圈车，或圈马，或把马和车一起圈成一个系统），然后\emph{只画外界对它的力}，别的什么都不画。边界一换，问题的面貌就换：把“马加车”圈成一个系统，马拉车、车拉马都成了系统内部的力，统统不用管，系统前进只凭地面向前推马蹄的摩擦力减去地面的阻力。画对受力图，力学题就完成了一半；画错受力图——比如把作用在别的物体上的力也塞进来——再熟练的计算也救不回来。

\step{一句话模型}

\begin{oneline}
力不决定你此刻跑多快，只决定你的速度正在变多快；而且力从来成双出现——你推世界多用力，世界就推你多用力，两个力分别落在两个物体上。
\end{oneline}

这句话其实把三条定律都装了进去：前半句是第一、第二定律（运动的默认状态是不变，合力只负责制造变化），后半句是第三定律。背下它，本章就不会白读。

\step{数学翻译器}

\begin{mainline}
把上一节的话翻译成式子，只需一行：
\[
  F_{\text{合}} = ma
\]
读法比写法重要。左边 \(F_{\text{合}}\) 是“外界对这个物体的所有推拉的加总”，回答的问题是：外界此刻总共想怎样改变它？右边的 \(a\) 是加速度，回答的问题是：它的速度实际上正在变多快？中间的 \(m\) 坐在分母的位置上（写成 \(a=F_{\text{合}}/m\) 更清楚）：质量越大，同样的合力能买来的加速度越小——这正是“质量是对加速的抗拒”的数学面孔。

新公式要先过三道例行检查。第一道，令 \(F_{\text{合}}=0\)，得到 \(a=0\)：速度不变——第一定律自动出现，说明三条规则是同一个模型，不是三条互不相干的标语。第二道，把 \(m\) 换成两倍，\(a\) 减半：装满大米的购物车起步慢一半，和生活经验一致。第三道，看方向：\(a\) 永远跟 \(F_{\text{合}}\) 同向，但完全可以跟速度\emph{反}向——刹车时合力向后、车仍向前走，速度在减小。旧模型最怕的这一幕，在新模型里只是再普通不过的一行算术。

再练一次带数量的感觉。一辆空购物车质量约 \(15\) 千克，你用 \(30\) 牛的力（大致相当于托起三瓶一升装可乐的力气）水平推它，忽略摩擦时加速度为 \(a=30/15=2\) 米每二次方秒：从静止推一秒，速度就到每秒 \(2\) 米，推两秒到每秒 \(4\) 米——难怪空车一碰就窜。装满大米后总质量变成 \(60\) 千克，同样的 \(30\) 牛只换来 \(0.5\) 米每二次方秒的加速度，四倍的“抗拒”，四分之一的成绩。这就是 \(a=F_{\text{合}}/m\) 在超市过道里的样子。
\end{mainline}

现在做一次带真实数据的估算，看看这条式子如何管着人命。一辆小汽车以约 \(50\) 千米每小时（约每秒 \(14\) 米）的速度撞墙停下，车里没系安全带的乘客几乎瞬间撞上仪表台，身体在约 \(0.05\) 秒内从每秒 \(14\) 米减到零。他的加速度大小约为
\[
  a = \frac{14\ \text{米/秒}}{0.05\ \text{秒}} = 280\ \text{米/秒}^2,
\]
将近重力加速度的三十倍。一个 \(60\) 千克的人要承受 \(F=ma\approx 60\times 280 \approx 17000\) 牛的力——相当于一辆小汽车的重量压在身上。而系上安全带、撞上安全气囊时，安全带会拉伸、气囊会泄气，把减速过程拉长到约 \(0.5\) 秒：加速度降到约每秒 \(28\) 米，力降到约 \(1700\) 牛，大约是三个成年人的体重压在胸前——疼，但通常能活。注意这里什么也没有“变小”：速度的变化量是一样的，\(F_{\text{合}}=ma\) 只是把同一笔账分摊到了更长的时间里。工程师设计安全带、气囊、自行车头盔和快递箱里的泡沫，靠的都是这一行式子。这个想法——“同样的速度变化，时间越长，力越小”——将在第 10 章长成动量定理，今天先记住它的样子。

\begin{mainline}
再看受力图在公式里的位置。桌上放着一本书，画受力图（下图）：地球向下拉书（重力），桌面向上推书（支持力），两个力都作用在\emph{书}上，方向相反、大小相等，合力为零，所以 \(a=0\)，书静止。这里有一个本章最容易踩的坑：这一对力\emph{不是}作用力与反作用力！它们作用在同一物体（书）上，而第三定律的那对力必须分属两个物体——重力的反作用力是“书向上拉地球”，支持力的反作用力是“书向下压桌面”。平衡（同一物体上合力为零）和第三定律（两个物体间的成对出现）是两件不同的事，碰巧常常同时出现而已。
\end{mainline}

\begin{center}
\begin{tikzpicture}[scale=1.0]
  % 桌面
  \draw[thick] (-2.6,0) -- (2.6,0);
  \fill[pattern=horizontal lines] (-2.6,-0.28) rectangle (2.6,0);
  % 书
  \draw[thick,fill=gray!15] (-0.9,0) rectangle (0.9,0.7);
  \node at (0,0.35) {书};
  % 两个力都画在书上
  \draw[-{Stealth[length=3mm]},very thick,blue!70!black] (0.25,0.35) -- (0.25,1.85);
  \node[right=2pt] at (0.25,1.85) {支持力（桌面推书）};
  \draw[-{Stealth[length=3mm]},very thick,red!70!black] (-0.25,0.35) -- (-0.25,-1.15);
  \node[right=2pt] at (-0.25,-1.15) {重力（地球拉书）};
\end{tikzpicture}
\end{center}

\begin{mathbridge}
力是有方向的量，所以第二定律的完整写法是矢量式：
\[
  \vec{F}_{\text{合}} = m\vec{a},
  \qquad\text{其中}\qquad
  \vec{F}_{\text{合}} = \vec{F}_1 + \vec{F}_2 + \cdots
\]
实际解题时沿互相垂直的两个方向拆开：水平方向写 \(F_{\text{合},x}=ma_x\)，竖直方向写 \(F_{\text{合},y}=ma_y\)，各算各的，互不干扰。桌上的书就是 \(F_{\text{合},y}=0\) 的例子：支持力与重力相消，竖直加速度为零。

牛顿本人写的其实不是 \(F=ma\)，而是更基本的形式
\[
  \vec{F}_{\text{合}} = \frac{d\vec{p}}{dt},
  \qquad \vec{p}=m\vec{v}\ \text{（动量）},
\]
即“合力等于动量随时间的变化率”。当质量不变时，\(\frac{d\vec{p}}{dt}=m\frac{d\vec{v}}{dt}=m\vec{a}\)，退化成我们熟悉的样子；但它还能处理质量会变的对象，比如一边喷气一边变轻的火箭。这一条式子的威力要在第 10 章才真正展开。

顺带回答一个也许你已经注意到的问题：本书开篇提出的三个追问——“系统此刻是什么状态？什么规律决定状态如何变化？我们怎样通过测量知道它？”——在牛顿力学里第一次有了完整的回答。状态由位置和速度给出；状态如何变化由 \(F_{\text{合}}=dp/dt\) 给出；而测量，就是用尺和钟去记录位置随时间的变化（加速度计和体重计量的其实都是力）。这个“状态—规律—测量”的三件套会贯穿全书，只不过后面每一篇都会换掉其中的零件。

最后看一眼图像语言：对同一个物体，\(a\) 与 \(F_{\text{合}}\) 成正比，图像是一条过原点的直线，斜率就是 \(1/m\)。质量越大，直线越“躺平”——同样的力，买回的加速度越少。
\end{mathbridge}

\begin{center}
\begin{tikzpicture}[scale=1.0]
  \draw[-{Stealth[length=3mm]}] (0,0) -- (5.0,0) node[below left] {合力 \(F_{\text{合}}\)};
  \draw[-{Stealth[length=3mm]}] (0,0) -- (0,3.0) node[below left] {加速度 \(a\)};
  \draw[very thick,blue!70!black] (0,0) -- (4.4,2.2) node[right=2pt] {\(m\) 小（车空着）};
  \draw[very thick,red!70!black] (0,0) -- (4.4,0.9) node[right=2pt] {\(m\) 大（装满大米）};
  \draw[dashed] (3,0) node[below] {\(F_0\)} -- (3,1.5) -- (0,1.5) node[left] {\(a_0\)};
\end{tikzpicture}
\end{center}

\begin{challenge}
火箭是第二定律最壮丽的应用，也顺带击碎了“火箭靠推地面或推空气上天”的误解。它把携带的燃料烧成燃气，高速向后喷出：火箭向后推燃气，燃气向前推火箭——第三定律在真空里照样成立，因为太空中根本不需要“有什么东西在后面顶着”。麻烦在于火箭的质量一直在减少：燃料烧掉一半，火箭就轻了一半，同样的推力买来的加速度越来越大。把 \(F=\frac{dp}{dt}\) 同时用于火箭和喷出的燃气，可以推出齐奥尔科夫斯基公式：火箭最终能获得的速度增量，取决于喷气速度和“出发时质量与烧完时质量之比”的对数。对数是个吝啬鬼——想再快一点，燃料要按倍数往上堆。这就是为什么运载火箭要造好几级，烧完一级扔一级。具体的推导留给第 10 章的挑战层，这里只需记住：太空里没有空气可推，火箭靠的是把一部分自己向后扔出去。
\end{challenge}

\step{模型的边界}

\begin{mainline}
牛顿三定律好用得近乎霸道，但它的领地有明确的国界，越界使用就会出错。先说适用条件，再说典型误用。

适用条件一：\emph{惯性参考系}。牛顿定律只在“不加速、不转动的参考系”里直接成立。地面（近似）可以，匀速直线行驶的车厢可以；但急刹车的车厢不行——你在车里“看到”人会凭空向前冲，没有任何水平力推他，这在车厢参考系里根本无法用牛顿定律解释。补救办法是引入“惯性力”，但那是一种为了在非惯性系里继续套用公式而发明的记账技巧，账本上必须注明它找不到施力物体。

适用条件二：\emph{速度远小于光速}。当物体接近光速，\(a=F/m\) 的预言开始系统性偏离实验：恒定推力下，速度的增长越来越慢，怎么也推不过光速。这不是公式“算错了一点”，而是速度、时间、质量之间的关系需要重写（第 38—40 章）。有趣的是，动量形式 \(F=dp/dt\) 比 \(F=ma\) 活得更久——相对论保留的正是“合力改变动量”这条内核。

适用条件三：\emph{宏观物体}。到了原子尺度，电子、原子这些对象不再沿着确定的轨道服从牛顿定律，需要用量子力学（第 43 章起）描述。这不是因为我们的仪器太粗糙测不准，而是“同时具有确定位置和确定速度”这个牛顿式的前提本身在微观世界不成立。
\end{mainline}

再列四种本章特有的典型误用，每一条都对应着前面埋过的雷：

\begin{itemize}[leftmargin=1.8em]
  \item \emph{“物体需要力来维持运动。”}错。合力改变的是动量，不是“维持”速度；匀速直线运动不需要任何力。这句话是两千年旧直觉的化石，本章整章都在拆它。
  \item \emph{把质量当重量。}质量是抗拒加速的属性，重量是地球对它的拉力。太空里重量没了，质量还在，推大箱子照样费劲。
  \item \emph{把作用力与反作用力加在同一个物体上然后宣布抵消。}马与马车悖论的翻版。记住：只有作用在同一物体上的力才谈得上加总和抵消。
  \item \emph{“速度方向在变的圆周运动，合力为零。”}错。速度是带方向的量，方向变也是变；拐弯中的汽车、旋转中的秋千，合力都不为零。这账怎么算，是下一章的主题。
\end{itemize}

\step{回头解释开场现象}

现在回到开头的三个悬念，用新模型逐一对账。

\textbf{冰壶。}松手之后，冰壶在水平方向只受冰面微小的摩擦力。合力很小，由 \(F_{\text{合}}=ma\)（或更基本地，合力决定动量变化的快慢），它的速度只能极缓慢地减小，于是滑出三十多米。草地上摩擦力大得多，同一笔“速度存款”被迅速扣光，半米就停。绝对光滑的冰面上合力为零——第一定律接管的区域——冰壶将以推出瞬间的速度匀速直线滑下去，永远不停。你在“先猜一猜”里押的（b）是对的：\emph{不停才是默认，停下才需要理由}。太空扳手是同一个故事的纯净版：周围连摩擦都没有，它便沿着直线永远漂下去。

\textbf{抽纸条。}慢拉时，纸与杯底之间的摩擦力作用了好几秒，杯子的速度被一点点带起来，于是跟着纸走。猛抽时，摩擦力并没有变大多少，但它只作用了零点几秒，杯子来不及积累多少速度，纸就已经溜走——“同样的推动，时间越短，造成的速度变化越小”，这正是安全带估算里那笔账的家用版。装满水的杯子质量更大，同样的摩擦推动换来的速度变化更小，所以更稳。你在家看到的，就是惯性坐在餐桌上的样子。

\textbf{公交刹车。}刹车时，车厢在减速，你的脚通过摩擦跟着车厢减速，可你的上半身没有受到向后的力，它忠实地执行第一定律，保持原来的速度继续前进——于是相对车厢向前倾。从头到尾没有“一股力把你往前甩”：在地面参考系里，你上半身只是“没被拉住”，继续匀速而已；那个“被甩”的感觉，是待在加速参考系（车厢）里产生的错觉。同理，启动时你向后仰，是脚被车厢带着向前、上半身落在了后面。安全带存在的意义，就是替你的上半身提供那个“向后拉”的力，让它和车厢一起减速，而不是替你甩出去。

三个悬念，同一把钥匙：把“运动”和“运动的变化”分开，合力只管后者。

\step{脑洞实验与挑战题}

\begin{thought}[无限冰面与宇宙尽头的冰壶]
把本章的想象推到极端尺度。假设有一块绝对光滑、无限大的冰面，你把冰壶轻轻一推，然后关掉全宇宙的灯，忘掉这件事。一千万年以后它在哪里？按第一定律：仍以当初那个速度，沿当初那条直线前进，走过的距离大约是“速度乘以一千万年”。再极端一点：一艘飞船飞到远离一切星系的深空，远到引力都微弱得可以忽略，然后关掉发动机。它会慢慢停下来吗？不会。它将以关机瞬间的速度永远匀速直线地漂下去，一亿年、一百亿年都一样——没有什么去消耗它的运动，因为运动本来就不需要被“维持”。这个思想实验反过来照出我们地球生活的不寻常：我们之所以觉得“停下来才自然”，是因为我们恰好住在一个到处充满摩擦和阻力的角落里。第一定律不是理想化的多余假设，它是被地球上的阻力层层盖住的世界默认值。
\end{thought}

\begin{thought}[掉落的电梯里，牛顿定律失灵了吗]
再开一个把本章连向远方理论的脑洞。设想你站在电梯里的体重计上，缆绳突然断了，电梯和你一起在重力作用下自由下落。这一瞬间你会发现：体重计读数变成零，松开手里的苹果，它不往下掉，就飘在你面前。在这个下落的电梯里，重力好像“消失”了。牛顿力学对此的解释是：电梯、你、苹果全都以同样的加速度下落，谁也不需要额外的力，所以你们之间没有相互挤压——重力还在，只是大家都被它同样地拉着走。但爱因斯坦盯上了另一种读法：在足够小的区域里，“自由下落的电梯”和“远离一切引力的深空”做不出任何力学实验能把二者区分开。既然区分不开，也许引力根本就不是一种普通的力，而是时空本身的几何性质——这条“等效原理”是广义相对论（第 41 章）的入口。本章的“质量是抗拒加速的属性”在这里也露出另一张脸：抗拒加速的惯性质量，和被地球吸引的引力质量，为什么偏偏永远相等？牛顿定律把它们当成碰巧相等，爱因斯坦把它们当成同一回事。一个藏在电梯里的疑问，最终掀翻了整个引力理论。
\end{thought}

\begin{puzzles}
\item 公交车向前行驶中急刹车，车厢里用细线悬着一只氦气球（它本来飘在车顶附近）。刹车瞬间，气球向车头飘还是向车尾飘？\emph{思路提示}：别先想气球，先想车里的空气。刹车时空气由于惯性向前涌，车头一侧气压略高、车尾略低。气球像一个“怕高压、爱低压”的东西——它总是从气压高处被挤向气压低处。想清楚空气往哪边堆，答案就出来了，而且和你的身体倾斜方向恰好相反。
\item 火箭在太空中飞行，周围是什么都没有的真空，没有地面可蹬、没有空气可推，它喷出的火焰“推”在什么上？火箭凭什么加速？\emph{思路提示}：把系统边界画大一点——把火箭和它已经喷出去的所有燃气圈成一个系统。火箭向后扔燃气，燃气反过来向前推火箭，第三定律在系统内部完成了全部工作。再想想：把船上带的沙袋一袋一袋使劲向后扔，船会怎样？
\item 你站在体重计上，突然快速下蹲。下蹲刚开始的一瞬间，体重计读数比你真实体重大、小、还是不变？\emph{思路提示}：刚开始下蹲时，你身体的重心在向下\emph{加速}。由 \(F_{\text{合}}=ma\)，向下的加速度要求合力向下，即重力大于支持力。体重计显示的是支持力。想清楚方向再回答；顺便预测一下，下蹲快要停住时读数又会怎样变化。
\item 生鸡蛋从同样的高度分别落在水泥地和厚海绵上，落在海绵上的通常不碎。两者停下前速度变化几乎一样，区别到底在哪？\emph{思路提示}：用安全带那笔账来想——“同样的速度变化，时间越短，力越大”。海绵把鸡蛋停下来的过程拉长了多少，蛋壳受到的力就小了多少。这也是快递箱里泡沫填充物的工作原理。
\end{puzzles}

至此，三条规则已经立起来了：运动默认不变（第一定律），合力决定动量如何变（第二定律），力永远成双、分属两个物体（第三定律）。但我们还没有认真清点过“力”这个家族本身的成员——重力、支持力、拉力、摩擦力、空气阻力，它们各自从哪来、服从什么规律、什么时候会“按需要变化”？那是下一章的任务。而“合力改变动量”这句本章反复念叨的话，将在第 10 章长成一整本关于碰撞的记账法。
